\chapter{Metodologie di analisi e tecnologie analizzate}
\label{cap:processi-metodologie}

\intro{Questo capitolo descrive le diverse fasi di ricerca ed analisi utilizzate per analizzare un prodotto a partire dallo studio di JasperServer mettendo in luce le diverse tecnologie considerate}

\section{Studio del prodotto JasperServer}
\label{sec:studio-jasperserver}
    La prima fase del progetto ha previsto uno studio approfondito del prodotto JasperServer, al fine di comprenderne le funzionalità, le caratteristiche tecniche e le modalità di utilizzo per poter .
    Per analizzare il prodotto JasperServer sono state utilizzate diverse metodologie di analisi, tra cui:
    \begin{itemize}
        \item \textbf{Analisi documentale}: è stata effettuata una revisione approfondita della documentazione ufficiale di JasperServer, inclusi manuali utente, guide di installazione e documenti tecnici. Questo ha permesso di comprendere le funzionalità principali del prodotto, le sue caratteristiche tecniche e le modalità di utilizzo.
        
        \item \textbf{Test pratici}: sono stati eseguiti test pratici utilizzando JasperServer in scenari reali o simulati. Questo ha permesso di valutare le prestazioni del prodotto, la facilità di configurazione e l'efficacia delle funzionalità offerte.
    \end{itemize}
    L'analisi di JasperServer ha fornito una base solida per la successiva fase di ricerca delle alternative, consentendo di identificare i requisiti chiave che le soluzioni alternative devono soddisfare, presenti alla tabella \ref{tab:requisiti-funzionali}.
\section{Prodotti analizzati}
\label{sec:prodotti-analizzati}
    Durante il progetto sono state analizzate diverse tecnologie con lo scopo di individuare quale di queste fosse la più adatta a sostituire JasperServer.
    Per eseguire questa ricerca sono stati consultati articoli, documentazioni ufficiali e tutorial online.
    \subsection{Birt}
        \emph{Birt} \footcite{site:birt} (Business Intelligence and Reporting Tools) è una piattaforma open source per la creazione di report e visualizzazioni di dati. È sviluppata dalla Eclipse Foundation ed è ampiamente utilizzata per lo sviluppo di applicazioni di business intelligence.
        Birt offre una vasta gamma di funzionalità, tra cui la creazione di report interattivi, la generazione di grafici e tabelle, l'integrazione con diverse fonti di dati e la possibilità di personalizzare l'aspetto dei report.
        La piattaforma Birt è basata su Java e può essere integrata in applicazioni web o desktop. Offre un ambiente di sviluppo intuitivo che consente agli sviluppatori di creare report complessi senza dover scrivere codice complesso.
        Birt supporta anche l'esportazione dei report in diversi formati, come PDF, Excel, HTML e Word, consentendo agli utenti di condividere facilmente i report generati.
        In sintesi, Birt è una piattaforma potente e flessibile per la creazione di report e visualizzazioni di dati, che offre una vasta gamma di funzionalità per soddisfare le esigenze delle applicazioni di business intelligence.
        \begin{figure}[H] 
            \centering 
            
\includegraphics[width=0.6\columnwidth]{tecnologie/BIRT-Logo.png} 
            \caption{Logo di BIRT}
        \end{figure}
    \subsection{Dynamic PDF}
        \emph{Dynamic PDF} \footcite{site:dynamicPDF}è una libreria software che consente la generazione dinamica di documenti PDF in modo programmatico. È sviluppata da Ceetron Solutions ed è ampiamente utilizzata per creare report, fatture, etichette e altri documenti PDF personalizzati.
        La libreria Dynamic PDF offre una vasta gamma di funzionalità, tra cui la creazione di pagine, l'aggiunta di testo, immagini, tabelle e grafici, nonché la formattazione avanzata dei contenuti. Supporta anche l'inserimento di collegamenti ipertestuali, segnalibri e moduli interattivi nei documenti PDF.
        Dynamic PDF è progettata per essere facile da usare e integrare in applicazioni software esistenti. Offre un'\emph{\gls{apig}} intuitiva che consente agli sviluppatori di creare documenti PDF in modo rapido ed efficiente, senza dover affrontare la complessità del formato PDF sottostante.
        In sintesi, Dynamic PDF è una libreria potente e flessibile per la generazione dinamica di documenti PDF, che offre una vasta gamma di funzionalità per soddisfare le esigenze delle applicazioni che richiedono la creazione di documenti PDF personalizzati.
        \begin{figure}[H] 
            \centering 
            \includegraphics[width=0.3\columnwidth]{tecnologie/DynamicPDF_Logo.png} 
            \caption{Logo di Dynamic PDF}
        \end{figure}
    \subsection{FineReport}
        \emph{FineReport} \footcite{site:fine-report} è una piattaforma di business intelligence e reportistica sviluppata da FanRuan Software Co., Ltd. È ampiamente utilizzata per la creazione di report, dashboard e analisi dei dati in ambito aziendale.
        FineReport offre una vasta gamma di funzionalità, tra cui la creazione di report interattivi, la generazione di grafici e tabelle, l'integrazione con diverse fonti di dati e la possibilità di personalizzare l'aspetto dei report. La piattaforma supporta anche l'esportazione dei report in diversi formati, come PDF, Excel, HTML e Word.
        FineReport è progettata per essere facile da usare e integrare in applicazioni software esistenti. Offre un ambiente di sviluppo intuitivo che consente agli sviluppatori di creare report complessi senza dover scrivere codice complesso.
        In sintesi, FineReport è una piattaforma potente e flessibile per la creazione di report e analisi dei dati, che offre una vasta gamma di funzionalità per soddisfare le esigenze delle applicazioni di business intelligence.
        \begin{figure}[H] 
            \centering 
            
\includegraphics[width=0.5\columnwidth]{tecnologie/FineReport-logo.png} 
            \caption{Logo di FineReport}
        \end{figure}    
    \subsection{Helical Insight}
        \emph{Helical Insight} \footcite{site:helical-insight} è una piattaforma di business intelligence open source che offre una vasta gamma di strumenti per l'analisi dei dati, la creazione di report e la visualizzazione delle informazioni aziendali. È progettata per essere flessibile e scalabile, consentendo alle organizzazioni di adattare la piattaforma alle proprie esigenze specifiche.
        Helical Insight offre funzionalità avanzate come l'integrazione con diverse fonti di dati, la creazione di dashboard interattivi, l'analisi predittiva e l'elaborazione in-memory. La piattaforma supporta anche l'esportazione dei report in diversi formati, come PDF, Excel, HTML e Word.
        Helical Insight è sviluppata in Java e può essere integrata in applicazioni web o desktop. Offre un ambiente di sviluppo intuitivo che consente agli sviluppatori di creare report complessi senza dover scrivere codice complesso.
        In sintesi, Helical Insight è una piattaforma completa di business intelligence che consente alle organizzazioni di sfruttare al massimo i propri dati per prendere decisioni informate e migliorare le prestazioni aziendali.
        \begin{figure}[H] 
            \centering 
            
\includegraphics[width=0.4\columnwidth]{tecnologie/Helical-insight-logo.png} 
            \caption{Logo di Helical Insight}
        \end{figure}
    \subsection{JsReport}
        \emph{JsReport} \footcite{site:jsreport} è una piattaforma di reportistica open source che consente la creazione, la gestione e la distribuzione di report in modo semplice ed efficiente. È progettata per essere flessibile e scalabile, consentendo agli sviluppatori di creare report personalizzati utilizzando tecnologie web standard come HTML, CSS e JavaScript.
        JsReport offre una vasta gamma di funzionalità, tra cui la generazione di report in diversi formati (PDF, Excel, HTML), l'integrazione con diverse fonti di dati, la creazione di dashboard interattivi e la possibilità di programmare l'esecuzione dei report. La piattaforma supporta anche l'uso di modelli personalizzati e l'estensione delle funzionalità tramite plugin.
        JsReport è sviluppata in Node.js e può essere integrata in applicazioni web o desktop. Offre un'\emph{\gls{apig}} intuitiva che consente agli sviluppatori di creare report complessi senza dover scrivere codice complesso.
        In sintesi, JsReport è una piattaforma potente e flessibile per la creazione di report, che offre una vasta gamma di funzionalità per soddisfare le esigenze delle applicazioni di reportistica.
        \begin{figure}[H] 
            \centering 
            \includegraphics[width=0.3\columnwidth]{tecnologie/JsReport_logo.png} 
            \caption{Logo di JsReport}
        \end{figure}
    \subsection{Knowage}
        \emph{Knowage} \footcite{site:knowage} è una piattaforma di business intelligence open source che offre una vasta gamma di strumenti per l'analisi dei dati, la creazione di report e la visualizzazione delle informazioni aziendali. È sviluppata da Engineering Ingegneria Informatica S.p.A. ed è ampiamente utilizzata da organizzazioni di diverse dimensioni per prendere decisioni basate sui dati.
        Knowage offre una suite completa di strumenti, tra cui Knowage Reporting per la creazione di report, Knowage OLAP per l'analisi multidimensionale dei dati e Knowage Dashboard per la creazione di dashboard interattivi. La piattaforma supporta l'integrazione con diverse fonti di dati, inclusi database relazionali, file di testo, servizi web e big data.
        Knowage offre anche funzionalità avanzate come l'analisi predittiva, l'elaborazione in-memory e l'integrazione con strumenti di visualizzazione dei dati come Tableau e QlikView.
        In sintesi, Knowage è una piattaforma completa di business intelligence che consente alle organizzazioni di sfruttare al massimo i propri dati per prendere decisioni informate e migliorare le prestazioni aziendali.
        \begin{figure}[H] 
            \centering 
            
\includegraphics[width=0.6\columnwidth]{tecnologie/KNOWAGE-logo.png} 
            \caption{Logo di Knowage}
        \end{figure}
    \subsection{Pentaho}
    \emph{Pentaho} \footcite{site:pentaho} è una piattaforma di business intelligence open source che offre una vasta gamma di strumenti
    per l'analisi dei dati, la creazione di report e la visualizzazione delle informazioni aziendali. È sviluppata da Hitachi Vantara ed è ampiamente utilizzata da organizzazioni di diverse dimensioni per prendere decisioni basate sui dati.
    Pentaho offre una suite completa di strumenti, tra cui Pentaho Data Integration (PDI) per l'estrazione, la trasformazione e il caricamento dei dati (ETL), Pentaho Business Analytics per l'analisi dei dati e la creazione di report, e Pentaho Dashboard Designer per la creazione di dashboard interattivi.
    La piattaforma supporta l'integrazione con diverse fonti di dati, inclusi database relazionali, file di testo, servizi web e big data. Offre anche funzionalità avanzate come l'analisi predittiva, l'elaborazione in-memory e l'integrazione con strumenti di visualizzazione dei dati come Tableau e QlikView.
    In sintesi, Pentaho è una piattaforma completa di business intelligence che consente alle organizzazioni di sfruttare al massimo i propri dati per prendere decisioni informate e migliorare le prestazioni aziendali.
    \begin{figure}[H] 
        \centering 
        
\includegraphics[width=0.2\columnwidth]{tecnologie/Pentaho-logo.png} 
        \caption{Logo di Pentaho}
    \end{figure}
    \subsection{ReportServer}
        \emph{ReportServer} \footcite{site:report-server} è una piattaforma di reportistica open source che consente la creazione, la gestione e la distribuzione di report in modo semplice ed efficiente. È progettata per essere flessibile e scalabile, consentendo agli sviluppatori di creare report personalizzati utilizzando tecnologie web standard come HTML, CSS e JavaScript.
        ReportServer offre una vasta gamma di funzionalità, tra cui la generazione di report in diversi formati (PDF, Excel, HTML), l'integrazione con diverse fonti di dati, la creazione di dashboard interattivi e la possibilità di programmare l'esecuzione dei report. La piattaforma supporta anche l'uso di modelli personalizzati e l'estensione delle funzionalità tramite plugin.
        ReportServer è sviluppata in Java e può essere integrata in applicazioni web o desktop. Offre un'\emph{\gls{apig}} intuitiva che consente agli sviluppatori di creare report complessi senza dover scrivere codice complesso.
        In sintesi, ReportServer è una piattaforma potente e flessibile per la creazione di report, che offre una vasta gamma di funzionalità per soddisfare le esigenze delle applicazioni di reportistica.
        \begin{figure}[H] 
            \centering 
            
\includegraphics[width=0.3\columnwidth]{tecnologie/ReportServerLogo.png} 
            \caption{Logo di ReportServer}
        \end{figure}
    \subsection{SAP Crystal Report}
        \emph{SAP Crystal Reports} \footcite{site:SAPcrystal} è uno strumento di reportistica sviluppato da SAP che consente agli utenti di creare report interattivi e visivamente accattivanti basati su dati provenienti da diverse fonti. È ampiamente utilizzato in ambito aziendale per generare report dettagliati e analisi dei dati.
        SAP Crystal Reports offre una vasta gamma di funzionalità, tra cui la creazione di report personalizzati, l'integrazione con diverse fonti di dati (database relazionali, file di testo, servizi web), la generazione di grafici e tabelle, e la possibilità di esportare i report in diversi formati (PDF, Excel, Word).
        Lo strumento supporta anche funzionalità avanzate come la formattazione condizionale, i sottoreport, le formule personalizzate e l'integrazione con applicazioni .NET e Java.
        In sintesi, SAP Crystal Reports è uno strumento potente e flessibile per la creazione di report aziendali, che offre una vasta gamma di funzionalità per soddisfare le esigenze delle organizzazioni che necessitano di analisi dei dati e reportistica avanzata.
        \begin{figure}[H] 
            \centering 
            
\includegraphics[width=0.3\columnwidth]{tecnologie/SAPCrystalReport_logo.png} 
            \caption{Logo di SAP Crystal Reports}
        \end{figure}
    \subsection{Stimulsoft}
        \emph{Stimulsoft} \footcite{site:stimulsoft} è una suite di strumenti di reportistica e business intelligence che consente agli utenti di creare report interattivi e dashboard personalizzati. È progettata per essere facile da usare e altamente flessibile, supportando diverse fonti di dati e offrendo un'ampia gamma di funzionalità per la generazione e la visualizzazione dei dati.
        Stimulsoft offre un ambiente visivo intuitivo per la creazione di report, con supporto per diversi formati di output come PDF, Excel, Word e HTML. Inoltre, consente l'integrazione con applicazioni web e desktop, fornendo \emph{\gls{apig}} e librerie che facilitano l'automazione della generazione dei report.
        In sintesi, Stimulsoft è uno strumento potente e versatile per la creazione di report aziendali, che offre una vasta gamma di funzionalità per soddisfare le esigenze delle organizzazioni che necessitano di analisi dei dati e reportistica avanzata.
        \begin{figure}[H] 
            \centering 
            
\includegraphics[width=0.5\columnwidth]{tecnologie/Stimulsoft_logo.png} 
            \caption{Logo di Stimulsoft}
        \end{figure}
\section{Strumenti utilizzati}
\label{sec:strumenti-utilizzati}
Per lo studio dei vari programmi di reportistica sono stati utilizzati i seguenti strumenti:
\begin{itemize}
    \item \emph{\textbf{Docker e Docker Compose}}: \footcite{site:docker}una piattaforma di containerizzazione che consente di creare, distribuire e eseguire applicazioni in ambienti isolati chiamati container. Docker semplifica la gestione delle dipendenze e garantisce la coerenza tra gli ambienti di sviluppo, test e produzione. Docker Compose è uno strumento che consente di definire e gestire applicazioni multi-container utilizzando un file di configurazione \texttt{docker-compose.yml}.
        In particolare, Docker è stato utilizzato per creare un ambiente di sviluppo isolato per le diverse applicazioni, consentendo di eseguire le applicazioni in un container separato dal sistema operativo host. Docker Compose è stato utilizzato per definire e gestire i servizi necessari per l'esecuzione degli applicativi, come il database e il server applicativo, dove possibile.
        \begin{figure}[H] 
            \centering 
            
\includegraphics[width=0.4\columnwidth]{strumenti/docker-logo.png} 
            \caption{Logo di Docker}
        \end{figure}
    \item \emph{\textbf{Postman}}: \footcite{site:postman}uno strumento utilizzato per testare le API RESTful. Consente di inviare richieste HTTP e visualizzare le risposte, facilitando il processo di sviluppo e debug delle \emph{\gls{apig}}. E' stato utilizzato per creare un server mock sul quale eseguire test sulle \emph{\gls{apig}} messe a disposizione dai vari software di reportistica e simulare passaggi di dati in formato JSON.
        \begin{figure}[H] 
            \centering 
            
\includegraphics[width=0.4\columnwidth]{strumenti/postman-logo.png}
            \caption{Logo di Postman}
        \end{figure}
    %\item \emph{\textbf{Eclipse IDE}}: \footcite{site:eclipse}un ambiente di sviluppo integrato (IDE) open source utilizzato per lo sviluppo di applicazioni Java. Offre funzionalità avanzate come il completamento automatico del codice, il debug e la gestione dei progetti.
    %    \begin{figure}[H]
    %        \centering 
    %        
\includegraphics[width=0.4\columnwidth]{strumenti/eclipse-logo.png} 
    %        \caption{Logo di Eclipse}
    %    \end{figure}
    
    %\item \emph{\textbf{Visual Studio Code}}: \footcite{site:visual-studio-code} un editor di codice sorgente leggero e potente, utilizzato per scrivere e modificare il codice del PoC. Supporta una vasta gamma di linguaggi di programmazione e offre funzionalità come il completamento automatico del codice, il debug e l'integrazione con Git.
    %    \begin{figure}[H] 
    %        \centering 
    %        
\includegraphics[width=0.2\columnwidth]{strumenti/vs-code-logo.png} 
    %        \caption{Logo di Visual Studio Code}
    %    \end{figure}
    %\item \emph{\textbf{VisualVM}}: \footcite{site:visualvm} uno strumento di monitoraggio e profilazione delle applicazioni Java. Consente di analizzare le prestazioni delle applicazioni, identificare i colli di bottiglia e ottimizzare l'utilizzo delle risorse.
    %    \begin{figure}[H] 
    %        \centering 
    %        
\includegraphics[width=0.2\columnwidth]{strumenti/visualvm-logo.png} 
    %        \caption{Logo di VisualVM}
    %    \end{figure}
\end{itemize}
\section{Analisi comparativa delle tecnologie}
\label{sec:analisi-comparativa-tecnologie}
La seguente fase del progetto ha previsto un'analisi comparativa delle tecnologie individuate nella sezione \ref{sec:prodotti-analizzati}, al fine di valutare le loro caratteristiche, funzionalità e prestazioni rispetto ai requisiti identificati nello studio di JasperServer, definiti nella tabella \ref{tab:requisiti-funzionali}.\\
Per ciascuna piattaforma sono state analizzate le informazioni disponibili nella documentazione ufficiale, includendo guide di installazione, manuali utente, riferimenti \emph{\gls{apig}} ed esempi di utilizzo. Questa fase ha consentito di acquisire una visione preliminare delle funzionalità offerte, dell'architettura del sistema e delle modalità di integrazione con applicazioni esterne.\\
Particolare attenzione è stata posta alla disponibilità di un editor grafico per la progettazione dei report, al supporto per sorgenti dati eterogenee (database \emph{\gls{sqlg}} e NoSQL, formati JSON) e alle modalità di generazione dei report tramite \emph{\gls{apig}}. Sono stati inoltre valutati gli aspetti relativi al deploy della piattaforma, verificando la possibilità di installazione su server locali e l'eventuale dipendenza da soluzioni cloud.\\
Lo studio della documentazione ha incluso anche una valutazione preliminare dell'esperienza di sviluppo, analizzando la chiarezza della documentazione, la presenza di esempi e la maturità delle \emph{\gls{apig}} fornite. Infine, sono stati considerati gli aspetti di licensing, al fine di individuare eventuali limitazioni o vincoli incompatibili con il contesto aziendale.\\
Sulla base delle informazioni raccolte, è stato possibile effettuare una prima selezione delle soluzioni più promettenti, individuando quelle che soddisfacevano i requisiti obbligatori e che risultavano maggiormente adatte a essere approfondite attraverso la realizzazione di un Proof of Concept.\\
Lo studio della documentazione ufficiale ha permesso di ridurre il rischio di implementazione, orientando la scelta verso soluzioni coerenti con i requisiti aziendali prima di procedere alla fase di sperimentazione pratica.\\
Di seguito viene riportata una tabella riassuntiva delle tecnologie analizzate, evidenziando la tipologia, i principali motivi di inclusione e i limiti principali rispetto ai requisiti individuati:
\begin{table}[H]%
\centering
\renewcommand{\arraystretch}{1.3}
\begin{tabularx}{\textwidth}{|>{\raggedright\arraybackslash}X
                              |>{\raggedright\arraybackslash}X
                              |>{\raggedright\arraybackslash}X
                              |>{\raggedright\arraybackslash}X|}
\hline
    \rowcolor{gray!40}
    \textbf{Soluzione} & \textbf{Tipologia} & \textbf{Motivo di inclusione} & \textbf{Limiti principali} \\
\hline
    BIRT + Knowage &
    Open-source &
    Combina editor grafico e server di gestione report; copre gran parte dei requisiti funzionali &
    Sottoreport non supportati in modo chiaro; supporto JSON solo tramite script manuali; upload report non automatizzato \\
\hline
    DynamicPDF &
    Commerciale &
    Buon supporto alla generazione PDF e ai sottoreport &
    Output limitato al PDF; assenza di \emph{\gls{apig}} di generazione; supporto JSON limitato \\
\hline
    FineReport &
    Commerciale &
    Editor grafico avanzato; server locale; diversi formati di output &
    Mancanza di \emph{\gls{apig}}; JSON solo locale e tramite scripting; assenza di sottoreport \\
\hline
    Helical Insight &
    Commerciale &
    Soluzione completa lato report e server; buona facilità d’uso &
    Versione open-source non più supportata; assenza di \emph{\gls{apig}}; sottoreport non disponibili \\
\hline
    JasperReport &
    Open-source &
    Buon supporto \emph{\gls{apig}}; JSON nativo; server locale; diversi formati di output &
    Assenza di editor grafico visuale; report da definire in HTML; sottoreport non supportati \\
\hline
    Pentaho Report Designer + Server &
    Commerciale &
    Ambiente di design maturo; supporto sottoreport; server locale &
    JSON solo tramite file locale; assenza di \emph{\gls{apig}} di generazione \\
\hline
    Report Server &
    Open-source e Commerciale &
    Piattaforma di gestione centralizzata dei report; supporto a diversi formati &
    Non include editor grafico; copre solo la parte server \\
\hline
    Stimulsoft Designer + Server &
    Commerciale &
    Piattaforma quasi completa; editor grafico avanzato; supporto JSON e sottoreport &
    \emph{\gls{apig}} non native; installazione server limitata (Windows / Docker) \\
\hline
    SAP Crystal Report &
    Commerciale &
    Editor grafico completo; supporto sottoreport; ampia diffusione &
    Solo ambiente Windows; JSON non nativo; non include server di gestione \\
\hline
\end{tabularx}
\caption{Confronto tra soluzioni di reportistica analizzate}
\label{tab:confronto-reporting}
\end{table}

Sulla base di questa analisi preliminare, alcune soluzioni sono state escluse in quanto non soddisfacevano i requisiti obbligatori definiti, mentre altre sono state ritenute idonee per un'analisi più approfondita. \\
In particolare le soluzioni escluse sono state le seguenti: DynamicPDF, FineReport, Helical Insight, Jsreport, Pentaho Report Designer + Server, SAP Crystal report.
Quelle sottoposte a un'analisi più approfondita sono Birt + Knowage, Report Server e Stimulsoft Designer + Server in quanto rispettano un maggior numero di requisiti obbligatori.

\section{Analisi approfondita delle tecnologie}
\label{sec:analisi-approfondita-tecnologie}
La fase successiva del progetto ha previsto un'analisi approfondita delle soluzioni selezionate, al fine di valutare in modo più dettagliato le loro funzionalità, prestazioni e facilità d'uso. \\
\subsection{Birt Design Report}
\label{sec:birt-design-report}
Il prodotto Birt Design Report è un editor per report realizzato sulla piattaforma Eclipse che permette di realizzare report tramite un editor grafico in modo similare a quanto offerto da JasperStudio ed è stato considerato per la parte di realizzazione di report in locale tramite editor grafico.
Viene offerto in forma open-source attualmente in versione 4.22.0 ma per l'analisi è stata considerata la versione 4.21.0 in quanto era quella disponibile in quel momento.\\
L'analisi nella sua prima fase si è concentrata nel verificare la presenza e la possibilità di utilizzare dati passati in formati JSON e forniti tramite parametro per popolare il report. 
Da questa analisi è emerso che per poter replicare questa funzionalità non presente nativamente è necessario definire uno script in Javascript che dato in input un parametro di tipo stringa lo utilizzi per scaricare i dati in formato JSON successivamente utilizzabili per riempire il report. 
Per realizzare questo si è proceduto in questo modo:
\begin{itemize}
    \item Definire un data source di tipo Scripted Data Source;
    \item Aggiungere un parametro, che chiameremo json, in Report Parameters;
    \item Creare un dataset collegandolo al datasource precedentemente definito;
    \item Sul dataset appena creato, nel file di design seleziono \texttt{Script} e successivamente in alto seleziono \texttt{open} dove aggiungo il seguente codice:

\end{itemize}

\subsection{Knowage}
\label{sec:knowage}
Questo prodotto è stato considerato per la parte server per replicare le funzionalità di JasperServer in quanto mette a disposizione un'implementazione grafica per ospitare report realizzati con Birt.
Per poterlo analizzare è stato scelto di utilizzare un'immagine docker ufficiale che fornisce il prodotto in versione 9.0, in quanto è una soluzione che permette di eseguire l'applicazione in un ambiente isolato e facilmente replicabile,
e anche l'installazione locale, che fornisce il prodotto in versione 8.1.7 in modo da poter verificare nel modo più completo possibile l'effettiva disponibilità delle funzionalità offerte.\\
L'analisi in questo punto è stata focalizzata principalmente sulla possibilità di eseguire il deploy di report da programma Birt Design Report, sulla possibilità di eseguire i report tramite \emph{\gls{apig}} e sulla possibilità di passare dati in formato JSON tramite parametro.\\
Le prove per poter eseguire il deploy dei report da Birt Design Report sono state effettuate sia tramite l'interfaccia grafica di Knowage che tramite \emph{\gls{apig}}.\\
Per quanto riguarda l'interfaccia grafica, essa presenta alcune discrepanze rispetto alla documentazione ufficiale, in quanto non è presente una sezione dedicata al deploy dei report ne al caricamento degli stessi, come è possibile verificare nelle immagini sottostanti.\\
\begin{figure}[H] 
    \centering
    \includegraphics[width=0.9\columnwidth]{analisi-tecnologie/Knowage_impostazioni_ufficiale.png}
    \caption{Interfaccia grafica di Knowage ufficiale}
\end{figure}
\begin{figure}[H] 
    \centering
    \includegraphics[width=0.9\columnwidth]{analisi-tecnologie/Knowage_impostazioni_locale.png}
    \caption{Interfaccia grafica di Knowage installata}
\end{figure}
Da successive verifiche non sono emerse altre versioni disponibili se non quelle già considerate, per cui è stato possibile escludere questa funzionalità.\\

\subsection{Knowage Studio}
\label{sec:knowage-studio}
Knowage Studio è un ambiente di sviluppo integrato (IDE) per la creazione e la gestione di report e dashboard all'interno della piattaforma Knowage. È progettato per semplificare il processo di sviluppo, consentendo agli utenti di creare report complessi e dashboard interattivi in modo intuitivo.\\
E' basato su Eclipse e offre un'interfaccia grafica che consente agli sviluppatori di progettare report utilizzando un editor drag-and-drop, senza la necessità di scrivere codice complesso. Knowage Studio supporta l'integrazione con diverse fonti di dati, inclusi database relazionali, file di testo e servizi web, consentendo agli utenti di accedere facilmente ai dati necessari per i loro report.
E' stato utilizzato da Knowage fino alla versione 6.2, poi definitivamente abbandonato a favore di BIRT.\\
Si è voluto analizzare questo prodotto in quanto poteva dimostrare l apossibilità di deploy di report in knowage tramite un editor grafico. Esso infatti funziona tramite l'aggiunta di un plugin a Eclipse che permette di caricare i report in modo simile a quanto fa JasperStudio per JasperServer.\\
L'analisi ha evidenziato che questo prodotto non è più supportato da Knowageper cui non è stato possibile testare questa funzionalità.

\subsection{Report Server}
\label{sec:report-server}
Report Server è una piattaforma di gestione centralizzata dei report che consente agli utenti di creare, gestire e distribuire report in modo semplice ed efficiente. È progettato per essere flessibile e scalabile, consentendo alle organizzazioni di adattare la piattaforma alle proprie esigenze specifiche.\\
Report Server è stato considerato per la parte server in quanto è un prodotto in grado di eseguire report sia jasper che Birt e quindi favorevole a un'implementazione più veloce e semplice come sostituto di JasperServer. Come nei prodotti precedenti è stata eseguita un'analisi delle funzionalità per verificarne il funzionamento e la possibilità di utilizzo.\\
Da questa analisi è emerso che Report Server mette a disposizione la possibilità di ospitare ed eseguire report realizzati sia in Jasper che in Birt ma non permette una loro esecuzione dinamica tramite script Javascript ma solo con uno script in linguaggio Groovy, che non è un linguaggio di scripting comune e potrebbe rappresentare una barriera per l'adozione da parte degli utilizzatori.\\
Essendo questo un limite importante e considerando il target di utenti a cui è rivolto il prodotto, si è deciso di escludere questa soluzione in quanto non soddisfa i requisiti obbligatori definiti in precedenza.\\

\subsection{Stimulsoft Designer e Stimulsoft Server}
\label{sec:Stimulsoft}
Stimulsoft è una suite di strumenti di reportistica e business intelligence che consente agli utenti di creare report interattivi e dashboard personalizzati. È progettata per essere facile da usare e altamente flessibile, supportando diverse fonti di dati e offrendo un'ampia gamma di funzionalità per la generazione e la visualizzazione dei dati.\\
Stimulsoft è stato considerato in combinazione tra Stimulsoft Designer e Stimulsoft Server per poter replicare le funzionalità di JasperServer, in quanto offre un editor grafico avanzato per la creazione dei report e un server per la gestione e l'esecuzione degli stessi.\\
L'ambiente di lavoro si presente in forma di un'applicazione desktop con grafica che ricorda molto quella di office, che quindi facilità l'utilizzo da parte di utenti non esperti.\\
Questo, pur essendo un grande punto di forza, rappresenta anche un limite in quanto il prodotto è disponibile solo per ambiente Windows e non è possibile eseguirlo su server Linux, che è l'ambiente di lavoro più comune per questo tipo di applicazioni.\\
Una modalità alternativa di esecuzione considerata è quella tramite Docker, ma in questo caso non si è riusciti pienamente ad eseguire il prodotto in quanto l'immagine docker ufficiale disponibile non  si eseguiva correttamente e quindi limitando la possibilità di testare le funzionalità offerte.\\

\subsection{Conclusioni}
A termine di questa analisi decisamente più approfondita dei vari prodotti e delle loro funzionalità, è stato possibile escludere quasi tutte le soluzioni considerate in quanto molte di esse non soddisfaceva completamente i requisiti obbligatori definiti in precedenza o non è sato possibile verificarne a fondo il funzionamento e quindi le fuzionalità offerte.\\
Il prodotto quindi tenuto a maggior considerazione per una successiva base di partenza per lo sviluppo di un PoC è stato Birt Design Report in quanto open-source in quanto permette una maggior personalizzazione, flessibiltà e facilità di integrazione con altre applicazioni, anche sviluppate internamente, che possono riuscire a replicare le funzionalità ricercate e definite.\\
Inoltre, Birt Design Report è un prodotto maturo e ampiamente utilizzato, con una vasta comunità di utenti e sviluppatori che contribuiscono al suo sviluppo e supporto. Questo garantisce una maggiore affidabilità e una maggiore disponibilità di risorse per la risoluzione di eventuali problemi o per l'implementazione di nuove funzionalità.\\

%\noindent Esempio di utilizzo di un termine nel glossario \\
%\gls{api}. \\

%\noindent Esempio di citazione in linea \\
%\cite{site:agile-manifesto}. \\

%\noindent Esempio di citazione nel pie' di pagina \\
%citazione\footcite{womak:lean-thinking} \\
